% Section 5.1, from lectures/L05/appendix/a_character_devices.md.

\section{Character Device Drivers}
\label{c5:app:a}

This section is how a driver presents itself to userspace: what a device node is, how the kernel
routes a \code{read()} on it into your function, and what each of those functions owes its caller.
\S\ref{c5:app:b} is the specification of the driver you write.

The framing worth keeping: \textbf{a character driver is a table of callbacks and nothing else.} You
fill in a \code{struct file_operations}, you register it against a number, and from then on the
kernel calls you. There is no main loop, no thread of yours, and nothing running between calls.

\subsection{The node, and the two numbers}
\label{c5:app:a1}

\S\ref{c2:app:b5} established that a device node carries a \textbf{major} and a \textbf{minor} where
a regular file has a size. Major selects the driver; minor selects which of that driver's devices.

Three facts follow, and all three surprise people the first time.

\textbf{A node is just a name.} \code{mknod} writes a directory entry containing a type and two
numbers, and consults nothing. You can create a node for a driver that does not exist; the failure
comes later, at \code{open}, as \code{ENXIO}.

\textbf{A node is not unique.} Two nodes with the same major and minor are the same device. Deleting
one removes a name and nothing else.

\textbf{The driver does not create the node.} Your driver registers a \emph{number}. Something else
makes a file appear in \file{/dev}. On this target that something is \textbf{devtmpfs}, prompted by
your call to \code{device_create}; on a desktop, \textbf{udev} sits on top and applies naming and
permission rules.

\subsection{Getting a number}
\label{c5:app:a2}

Two ways, and you want the second.

\begin{ccode}
/* Static: you pick the major and hope nobody else has it. */
register_chrdev_region(MKDEV(240, 0), 1, "qa_fifo");

/* Dynamic: the kernel picks, and tells you. */
alloc_chrdev_region(&devno, 0, 1, "qa_fifo");
\end{ccode}

\noindent Static allocation is how it was done for years and is a bad idea now: the major space is
shared, the ``local/experimental'' range is small, and a collision produces \code{-EBUSY} at load
time on somebody else's machine and not on yours. \textbf{Always allocate dynamically}, then find
out what you got:

\begin{console}
# grep qa_fifo /proc/devices
243 qa_fifo
\end{console}

\noindent 243 is not a constant. It is whatever was free, it can differ between boots, and any
script that hardcodes it is wrong. This is precisely why \code{device_create} and udev exist:
userspace should find the device by \emph{name}, and the number should be nobody's business.

\subsection{The table}
\label{c5:app:a3}

\begin{ccode}
static const struct file_operations qa_fops = {
        .owner   = THIS_MODULE,
        .open    = qa_open,
        .read    = qa_read,
        .write   = qa_write,
};
\end{ccode}

\noindent \textbf{Every member is optional except \code{.owner}.} A missing \code{.read} is not an
error; the kernel substitutes a default, and for \code{read} the default returns \code{-EINVAL}. So
a driver that forgets to implement something fails at the point userspace tries to use it, with an
error that says nothing about what is missing.

\textbf{\code{.owner = THIS_MODULE} is what stops the module being unloaded while a process has the
device open.} It is one line, it is easy to leave out, and leaving it out means \code{rmmod}
succeeds while a process still holds a file whose operations have just been freed. The result is not
an error message.

Registering the table against the number takes two more calls:

\begin{ccode}
cdev_init(&cdev, &qa_fops);
cdev.owner = THIS_MODULE;
cdev_add(&cdev, devno, 1);
\end{ccode}

\noindent After \code{cdev_add} the device is \textbf{live}: a \code{read()} on a matching node
reaches your function, even though you have not finished your init function yet. Everything your
callbacks depend on has to be ready before this call, not after it. That ordering constraint is the
single most common source of a crash-on-load in a first driver.

\subsection{The functions, and what they owe}
\label{c5:app:a4}

\subsubsection{\code{open} and \code{release}}
\label{c5:app:a:open-and-release}

\begin{ccode}
static int qa_open(struct inode *inode, struct file *file);
static int qa_release(struct inode *inode, struct file *file);
\end{ccode}

\noindent \code{open} returns 0 or a negative error. \code{release} is called when the \textbf{last}
descriptor referring to the open file is closed, which is not once per \code{close()}: a \code{fork}
duplicates descriptors, and \code{release} runs when the last of them goes.

\code{inode} identifies the device; \code{file} identifies this particular opening of it, and its
\code{private_data} is where per-open state belongs.

\subsubsection{\code{read} and \code{write}}
\label{c5:app:a:read-and-write}

\begin{ccode}
static ssize_t qa_read(struct file *file, char __user *buf, size_t count, loff_t *pos);
static ssize_t qa_write(struct file *file, const char __user *buf, size_t count, loff_t *pos);
\end{ccode}

\noindent The return value \textbf{is} the interface, and there are four cases:

\begin{booktable}{@{}>{\hsize=0.39\hsize}L>{\hsize=1.61\hsize}L@{}}
  \thead{Return} & \thead{Means} \\
  \midrule
  \code{n}, \code{0 < n <= count} & This many bytes were transferred. A short count is normal, not an error \\
  \code{0} from \code{read} & \textbf{End of file.} Nothing more will ever come \\
  \code{0} from \code{write} & Nothing was written, and it is not an error. Rarely what you mean \\
  negative & An error; the value is \code{-Exxx} and userspace sees it as \code{errno} \\
\end{booktable}

\noindent \textbf{The \code{0} row is the one this chapter exists for.} If your \code{read} returns
0 because it has nothing \emph{right now}, every reader concludes the device is finished. \code{cat}
exits, a loop terminates, and a program waiting for data quietly stops waiting. Nothing reports an
error, because nothing went wrong as far as userspace can tell.

Until Chapter~\ref{c9:ch} gives you wait queues, the honest answer when a driver has nothing to give
is \textbf{\code{-EAGAIN}}: ``not now, ask again''. Chapter~\ref{c9:ch} replaces it with blocking,
and \code{-EAGAIN} then becomes the answer only for \code{O_NONBLOCK} openers.

\subsubsection{Which error code}
\label{c5:app:a:which-error-code}

The code you return is part of the interface, because userspace acts on it.

\begin{booktable}{@{}lL@{}}
  \thead{Code} & \thead{Use for} \\
  \midrule
  \code{-EAGAIN} & No data now; try again \\
  \code{-ENOSPC} & No room; the write cannot be satisfied \\
  \code{-EFAULT} & A userspace pointer was bad. Return this when \code{copy_*_user} fails \\
  \code{-EINVAL} & The arguments do not make sense \\
  \code{-ENOMEM} & An allocation failed \\
  \code{-EIO} & The hardware failed \\
  \code{-ENODEV} & The device is gone \\
  \code{-ERESTARTSYS} & A signal arrived while sleeping. Chapter~\ref{c9:ch} \\
\end{booktable}

\lecturefigure{lectures/L05/appendix/images/read_contract.png}
  {Four situations a read faces, each with an arrow to what it must return: data available gives a byte
   count that may be short, empty and blocking sleeps, empty with O\_NONBLOCK gives -EAGAIN, and a
   signal gives -ERESTARTSYS. A red box below states that zero is not on the list, because it means
   end of file.}
  {fig:read_contract}

\subsection{The user pointer, which is not a pointer you may follow}
\label{c5:app:a5}

\code{buf} is an address \textbf{in the calling process's address space}. It looks like a \code{char
*} and the compiler will let you dereference it. Do not.

\begin{ccode}
memcpy(kbuf, buf, count);              /* WRONG */
if (copy_from_user(kbuf, buf, count))  /* right */
        return -EFAULT;
\end{ccode}

\noindent Three separate things are wrong with the first line, and it is worth having all three
rather than a rule to memorise:
\begin{itemize}
  \item \textbf{The page may not be present.} Userspace memory is demand-paged and can be swapped or
    not yet faulted in. A direct dereference from kernel context faults where no handler expects it.
  \item \textbf{The address may not belong to the caller at all.} It may point into kernel memory,
    and a driver that copies from it hands kernel memory to a userspace program that asked for it.
    That is not a theoretical class of bug; it is a large fraction of real kernel CVEs.
  \item \textbf{The mapping can change under you.} Another thread in the same process can
    \code{munmap} it between your check and your copy.
\end{itemize}
\code{copy_to_user} and \code{copy_from_user} check the range, handle the fault, and \textbf{return
the number of bytes they could not copy}, so zero means success, which reads backwards the first
time:

\begin{ccode}
if (copy_to_user(buf, kbuf, n))
        return -EFAULT;
\end{ccode}

\noindent They may also \textbf{sleep}, because faulting in a page can block. That is fine here and
is not fine in an interrupt handler, which is Chapter~\ref{c8:ch}.

\subsection{Making a node appear}
\label{c5:app:a6}

\code{cdev_add} makes the device reachable. It does not put anything in \file{/dev}.

\begin{ccode}
cls = class_create("qa_fifo");
device_create(cls, NULL, devno, NULL, "qa_fifo");
\end{ccode}

\noindent \code{class_create} makes a directory under \file{/sys/class/}, and \code{device_create}
adds a device to it, which emits a uevent that devtmpfs (or udev) turns into \file{/dev/qa_fifo}.

\textbf{Both of these have changed signature within recent memory}, which matters because most
tutorials you will find are older than the change:
\begin{itemize}
  \item \code{class_create} took an owner argument until 6.4 and takes only a name now.
  \item \code{no_llseek} was a real function for decades and \textbf{does not exist in 6.12}; code
    using it will not compile.
\end{itemize}
This is \S\ref{c3:app:a2}'s ``the internal API has no stability guarantee'' arriving in the first
driver you write. When a tutorial does not compile, the tutorial is usually right and old.

\subsection{Unwinding, in reverse}
\label{c5:app:a7}

Every one of those calls can fail, and each that succeeded must be undone in reverse order:

\begin{ccode}
        ...
        cls = class_create("qa_fifo");
        if (IS_ERR(cls)) { ret = PTR_ERR(cls); goto err_cdev; }
        ...
err_class:  class_destroy(cls);
err_cdev:   cdev_del(&cdev);
err_region: unregister_chrdev_region(devno, 1);
err_fifo:   qa_fifo_destroy(fifo);
        return ret;
\end{ccode}

\noindent The \code{goto} ladder is the kernel's house idiom for this and is not the sin it would be
elsewhere: each label undoes exactly one thing, and the order falls out of the order they were
acquired. Write it once here, because Chapter~\ref{c6:ch} shows you \code{devm_} and you delete the
whole ladder.

\code{IS_ERR} and \code{PTR_ERR} are the other idiom: functions returning pointers pack an error
into the pointer rather than using a separate output parameter. \code{IS_ERR(p)} tests it and
\code{PTR_ERR(p)} extracts the \code{-Exxx}. Checking for \code{NULL} instead is a bug that survives
testing, because the failure it misses is the rare one.

\subsection{The shortcut}
\label{c5:app:a8}

Everything above is about forty lines. The misc device is six:

\begin{ccode}
static struct miscdevice qa_misc = {
        .minor = MISC_DYNAMIC_MINOR,
        .name  = "qa_fifo",
        .fops  = &qa_fops,
};
misc_register(&qa_misc);
\end{ccode}

\noindent One shared major (10), a dynamically allocated minor, and the node created for you.
\textbf{Most small drivers should use it.} It is taught second because it hides the number, the
\code{cdev}, the class and the unwind path, and those are the parts worth having seen once.

\subsection{\code{ioctl}, and why to think twice}
\label{c5:app:a9}

\begin{ccode}
long qa_ioctl(struct file *file, unsigned int cmd, unsigned long arg);
\end{ccode}

\noindent \code{cmd} is an encoded number built with \code{_IO}, \code{_IOR}, \code{_IOW} or
\code{_IOWR}, packing a direction, a type, a number and the argument's size, so that a wrong call is
usually caught rather than misinterpreted.

The argument against adding one:
\begin{itemize}
  \item \textbf{It is an ABI you now own forever.} The moment a second program uses it, the numbers
    and the structure layout are frozen.
  \item \textbf{It is a portability trap.} A 32-bit process on a 64-bit kernel passes a differently
    sized structure, and handling that needs \code{compat_ioctl} and a struct with no implicit
    padding. Getting this wrong is a standing source of bugs.
  \item \textbf{There is usually a subsystem that already defines the operation.} That argument is
    Chapter~\ref{c11:ch}, and it is the reason this course's driver eventually deletes its character
    device entirely.
\end{itemize}
Use one when the operation genuinely is device-specific and no framework covers it. Reach for sysfs
for single values and debugfs for things that are not an ABI at all.
